\documentclass[conference]{IEEEtran}
\usepackage{times}
\usepackage[numbers]{natbib}
\usepackage{booktabs}
\usepackage{graphicx}
\usepackage{multirow}
\usepackage[bookmarks=true]{hyperref}

\pdfinfo{
   /Author (Anonymous)
   /Title  (Agentic Recovery for Lightweight Vision-Language-Action Policies)
   /CreationDate (D:20260606140000)
   /Subject (Robotics)
   /Keywords (robot learning; vision-language-action; failure recovery; semantic reasoning)
}

\begin{document}

\title{Agentic Recovery for Lightweight Vision-Language-Action Policies}

% Replace [add your ID here] after OpenReview assigns the SemRob submission ID.
\author{Author Names Omitted for Anonymous Review. Paper-ID [add your ID here]}

\maketitle

\begin{abstract}
% Author-written abstract only. RSS-style policy prohibits submitting text
% generated entirely by LLMs. Keep this concise: problem, method, evaluation,
% strongest quantitative result, and why SemRob should care.
TODO.
\end{abstract}

\IEEEpeerreviewmaketitle

\section{Introduction}
% Target: 4-page workshop version. Suggested shape:
% 1 paragraph: lightweight VLA policies are now practical but brittle.
% 1 paragraph: long-horizon failures often need detection and recovery.
% 1 paragraph: this work studies an external agentic recovery wrapper.
% 1 short contribution list.

TODO.

\section{Related Work}
% Keep this tight for SemRob. Three clusters are enough:
% efficient VLA models, agentic/planning robotics, and evaluation/recovery.

\subsection{Efficient Vision-Language-Action Policies}
TODO cite SmolVLA~\cite{shukor2025smolvla}, TinyVLA~\cite{wen2024tinyvla},
OpenVLA~\cite{kim2024openvla}, and broader robot foundation policies such as
RT-1~\cite{brohan2022rt1} and RT-2~\cite{zitkovich2023rt2}.

\subsection{Agentic Planning and Feedback}
TODO cite SayCan~\cite{ahn2022saycan}, Inner Monologue~\cite{huang2022inner},
Code as Policies~\cite{liang2022code}, ReAct~\cite{yao2022react}, and
Reflexion~\cite{shinn2023reflexion}. Tie this to SemRob topics: semantic
goal interpretation, policy steering, online correction, and uncertainty or
failure monitoring.

\subsection{Benchmarks and Failure-Recovery Evaluation}
TODO cite LIBERO~\cite{liu2023libero}. State why final success must remain the
benchmark success flag rather than a wrapper-internal verifier decision.

\section{Method}
% Core claim: no base policy weight changes. The wrapper controls when and how
% a fixed lightweight VLA is retried or replanned.

\subsection{Base Policy and Evaluation Harness}
TODO describe the fixed SmolVLA policy, observation/action interface, LIBERO
or ManiSkill environment, seeds, action budgets, and logging artifacts.

\subsection{Agentic Recovery Wrapper}
TODO define planner, verifier, retry scheduler, and replan policy. Distinguish
episode-level retry from subgoal-level in-episode recovery if both appear.

\subsection{Success Semantics}
TODO state explicitly: final task success is the environment benchmark success
signal, e.g., \texttt{info["success"]} or \texttt{is\_success}. Wrapper
verifier predicates are used only for control decisions such as continue,
retry, or replan.

\section{Experiments}
% Keep this evidence-first. Use matched seeds, same base policy, same action
% budget where possible. Do not overclaim from small probes.

\subsection{Setup}
TODO list suites, task IDs, episodes per task, compute environment, simulator
versions, policy checkpoint, and action-step protocol.

\subsection{Policy-Only Baseline}
TODO summarize the current routed SmolVLA baseline and cite the local artifact
paths in the camera-ready or supplement, not in the blinded submission.

\begin{table}[t]
\centering
\caption{Policy-only baseline reproduction on LIBERO. Fill with final
paper-scale numbers only.}
\label{tab:baseline}
\begin{tabular}{lrrrrr}
\toprule
Condition & Goal & Object & Spatial & Long & Avg \\
\midrule
Policy-only & TODO & TODO & TODO & TODO & TODO \\
Reference & TODO & TODO & TODO & TODO & TODO \\
\bottomrule
\end{tabular}
\end{table}

\subsection{Recovery Results}
TODO compare \texttt{policy\_only}, \texttt{agentic\_retry}, and, if ready,
\texttt{agentic\_replan}. Report recovery-after-failure, retries per episode,
latency, and failure categories.

\begin{table}[t]
\centering
\caption{Agentic recovery ablation. Keep all conditions matched by seeds,
base policy, and action budget unless explicitly stated.}
\label{tab:recovery}
\begin{tabular}{lrrrr}
\toprule
Condition & Episodes & Success & Recovery & Retries \\
\midrule
Policy-only & TODO & TODO & -- & -- \\
Agentic retry & TODO & TODO & TODO & TODO \\
Agentic replan & TODO & TODO & TODO & TODO \\
\bottomrule
\end{tabular}
\end{table}

\section{Discussion}
TODO explain when recovery helps, when it fails, and what this says about
semantic reasoning around lightweight VLA policies.

\section{Limitations}
% RSS requires a limitations section for main papers; SemRob follows RSS paper
% policy, so include it here too.
TODO include limited task coverage, simulator-only evidence, verifier
dependence, no weight updates, possible extra action budget, and no real robot
claim unless real evidence is added.

\section{Conclusion}
TODO one paragraph: what was tested, strongest measured result, and why the
wrapper is useful as an evaluation and recovery layer for lightweight VLA
policies.

\bibliographystyle{plainnat}
\bibliography{references}

\end{document}
